\DiaryEntry{Complex Analysis, 2}{2026-02-05}{Complex Analysis}

\paragraph{Complex Exponential.} The exponential function is defined as below,

\be\label{2026-02-05:eq1}
e^{z} = e^{x + jy} = e^x e^{jy} = e^x (\cos(y) + j \sin(y))
\ee

Let's consider some properties of the complex exponential. We have

\bee
e^z = 1 \, \text{iff} \, z = 2k\pi j
\eee

for any integer $k$. Using $z = x + jy$ we have $|e^z| = |e^x| |e^{jy}| = e^x = 1$ and therefore $x = 0$. So we have $e^z = e^{jy} = \cos y + j \sin y$. Equating real and imaginary parts, we have $\cos y = 1, \sin y = 0$ and this is satisfied iff $y = 2k\pi$ with $k$ being some integer.
In a similar spirit, we have

\bee
e^{z_1} = e^{z_2} \, \text{iff} \, z_1 = z_2 + 2k\pi j
\eee

Since, for all $z$, $e^{z+2k\pi j} = e^z$, we find that $e^z$ is periodic with complex period $2\pi j$. Consequently, if we divide the $z$-plane into the infinite horizontal strips

\bee
S_n = \{ x + jy: -\infty < x < \infty, \, (2n - 1) \pi < y \leq (2n+1) \pi, \, n = 0, \pm 1, \pm 2, \ldots \}
\eee

then $e^z$ will behave in the same manner on each strip. We also see that $e^z$ is one-to-one on each sttrip $S_n$.

We can insert \eqref{2026-02-05:eq1} into the Cauchy-Riemann equations to check whether the complex exponential is analytic,

\bee
\frac{\partial u}{\partial x} = \frac{\partial v}{\partial y}, \quad \frac{\partial u}{\partial y} = - \frac{\partial v}{\partial x}
\eee

We have

\begin{align*}
    \frac{\partial u}{\partial x} &= e^x \cos y, \quad \frac{\partial v}{\partial y} = e^x \cos y \\
    \frac{\partial u}{\partial y} &= - e^x \sin y, \quad \frac{\partial v}{\partial x} = e^x \sin y
\end{align*}

and see that the Cauchy-Riemann equations are fulfilled and conclude that the complex exponential is an analytic function. We can now obtain the complex derivative as

\bee
f'(z) = \frac{\partial u}{\partial x} + \frac{\partial v}{\partial x} = e^x \cos y + j e^x \sin y = e^z \qed
\eee


\paragraph{Trigonometric Functions.} We start with

\begin{align*}
    e^{j \theta} &= \cos \theta + j \sin \theta \\
    e^{-j \theta} &= \cos \theta - j \sin \theta
\end{align*}

Adding the two equations yields

\bee
\cos \theta = \frac{1}{2} \left( e^{j \theta} + e^{-j \theta} \right)
\eee

Subtracting yields

\bee
\sin \theta = \frac{1}{2j} \left( e^{j \theta} - e^{-j \theta} \right)
\eee

This inspires the definition of the trigonometric functions with complex argument as

\bee
\sin z= \frac{1}{2j} \left( e^{j z} - e^{-j z} \right), \,  \cos z = \frac{1}{2} \left( e^{j z} + e^{-j z} \right)
\eee

We can use these definitions (together with the fact that $e^z$ is analyti), to obtain first derivatives; eg

\begin{align*}
\frac{d \sin z}{dz} &= \frac{d }{dz} \frac{1}{2j} \left( e^{j z} - e^{-j z} \right) = \frac{1}{2j} \left( j e^{j z} + j e^{-j z} \right) \\
&= \frac{1}{2} \left( e^{j z} + e^{-j z} \right) = \cos z \qed
\end{align*}

We can use \eqref{2026-02-05:eq1} to (re)derive some trigonometric identities. We have

\bee
e^{j (z_1 + z_2)} = \cos (z_1 + z_2) + j \sin (z_1 + z_2)
\eee

However, we can also expand the exponential and then split into real and imaginary parts,

\begin{align*}
e^{j (z_1 + z_2)} &= e^{j z_1} \cdot e^{j z_2} = (\cos z_1 + j \sin z_1) \cdot (\cos z_2 + j \sin z_2) \\
&= \cos z_1 \cos z_2 - \sin z_1 \sin z_2 + j \cos z_1 \sin z_2 + j \cos z_2 \sin z_1
\end{align*}

Comparing real and imaginary parts, we obtain

\begin{align*}
    \cos (z_1 + z_2) &= \cos z_1 \cos z_2 - \sin z_1 \sin z_2 \\
    \sin (z_1 + z_2) &= \cos z_1 \sin z_2 + \cos z_2 \sin z_1
\end{align*}

These identities can be used to derive all sorts of interesting results; Equating

\bee
\sin(z + 2 \pi) = \cos z \sin 2 \pi + \cos 2 \pi \sin z = \sin z
\eee

Further hrigonometric functions, hyperbolic functions

\paragraph{Complex Logarithm.}


%%% Local Variables:
%%% mode: latex
%%% TeX-master: "journal"
%%% End:
